\section{A leakage bound for the original prolate trial}
\label{sec:spectral-prolate-rayleigh}

This section addresses the size of the Rayleigh quotient of the
specific two-mode trial in S25, equations \texttt{wop} and
\texttt{ktoh}.  It does not identify that trial with the lowest
eigenvector.  In particular, the estimate below is an upper bound,
not a two-sided eigenvalue asymptotic.

\begin{theorem}[AC-S30: prolate leakage controls the Weil Rayleigh quotient]
\label{AC-S30}
Use the Fourier transform
$\mathcal Ff(v)=\int_{\mathbb R}f(u)e^{-2\pi iuv}\,du$,
and let $P_\lambda=1_{[-\lambda,\lambda]}$.
Let $\phi_{0,\lambda},\phi_{4,\lambda}$ be real, even,
unit-norm, time-supported prolate functions, with the usual
Hermite indexing.  Write
\[
 P_\lambda\mathcal F\phi_{j,\lambda}
       =\mu_j\phi_{j,\lambda},\qquad
 \nu_j=\mu_j^2,\qquad d_j=\phi_{j,\lambda}(0),
             \quad j=0,4.                         \tag{SR.1}
\]
For all sufficiently large $\lambda$, $\mu_0,\mu_4>0$,
$d_0d_4\ne0$, and define
\[
 \begin{split}
 \psi_\lambda&=-\frac34\left(
          \mu_4\frac{\phi_{0,\lambda}}{d_0}
             -\mu_0\frac{\phi_{4,\lambda}}{d_4}\right),\\
 k_\lambda(x)&=1_{[-L,L]}(x)e^{x/2}
                 \sum_{n\geq1}\psi_\lambda(ne^x),
                  \qquad L=\log\lambda,\qquad
 \delta_\lambda=1-\nu_4.                           \end{split}
 \tag{SR.2}
\]
This is precisely the mean-zero two-mode trial of S25,
with a scalar chosen so that its limit is the $K$ of AC-S04.
It satisfies $k_\lambda\to K$ in $L^2(\mathbb R)$.

Under $A$, for every $0<\varepsilon<1$ there are
$C_\varepsilon,A_\varepsilon<\infty$ such that
\[
 q_L(k_\lambda)\leq
       C_\varepsilon\lambda^{A_\varepsilon}
                       \delta_\lambda^{1-\varepsilon}.
                                                        \tag{SR.3}
\]
In particular, for every $\eta>0$,
\[
 0<a_1(\log\lambda)
       \leq C_\eta\exp\bigl(-(4\pi-\eta)\lambda^2\bigr),
 \qquad
 \limsup_{\lambda\to\infty}
            \lambda^{-2}\log a_1(\log\lambda)\leq-4\pi.
                                                        \tag{SR.4}
\]
The same Rayleigh upper bound is available in the even
logarithmic sector, by replacing $k_\lambda(x)$ by
$(k_\lambda(x)+k_\lambda(-x))/2$.
Neither conclusion asserts that the unrestricted minimum
is simple or even.
\end{theorem}

\paragraph{The classical prolate input, with its normalization.}
The uniform comparison used below is S25, Lemma
\texttt{meixnerlem}, equations \texttt{esti01},
\texttt{meix0}--\texttt{esti}: after scalar normalization,
the modes $0,4$ differ uniformly on $[-\lambda,\lambda]$
from their Hermite limits by $O(\lambda^{-2})$.
Ratios by their nonzero values at zero have the same
estimate, independently of the $L^2$ normalization.

For the scalar leakage we use W.~H.~J. Fuchs,
\emph{On the eigenvalues of an integral equation arising
in the theory of band-limited signals},
J. Math. Anal. Appl. \textbf{9} (1964), 317--330,
doi:10.1016/0022-247X(64)90017-4, Theorem~1 on p.~319.
His equation (1), p.~317, has kernel
$\sin(x-y)/(\pi(x-y))$ on $[-c,c]$.  His parameter
$a$ satisfies $a^2=c$, and his eigenvalue is the
concentration eigenvalue, not its Fourier square root.
Rescaling our concentration kernel
$\sin(2\pi\lambda(u-v))/(\pi(u-v))$ gives
\[
 c=2\pi\lambda^2,\qquad
 1-\nu_n\sim\frac{4\sqrt\pi\,8^n}{n!}
                  c^{n+1/2}e^{-2c}.                \tag{SR.5}
\]
Thus
\[
 \delta_\lambda\sim
       \frac{2^{15}\sqrt2\,\pi^5}{3}
                         \lambda^9e^{-4\pi\lambda^2},
 \qquad
 1-\mu_4\sim\frac12\delta_\lambda.                \tag{SR.6}
\]
The second constant is exactly the one quoted in S25
after \texttt{meixnerlem}(ii).
Fuchs's proof runs through Lemmas 1--5, pp.~321--329,
and Section III, p.~329; it differentiates the eigenvalue
using the endpoint value, obtains that endpoint value by
matching the imaginary-axis Fourier and differential-equation
asymptotics, and integrates the resulting derivative.
The accessible original-article OCR is at
\url{https://www.scribd.com/document/466874374/1-s2-0-0022247X64900174-main}.
The audit read its complete accessible proof text; some
intermediate displayed WKB formulas have defective OCR,
so this is recorded as use of a classical theorem with
an access limitation, not as a complete independent
verification of every formula in that proof.  The
new estimate (SR.3) is stated in the exact leakage
$\delta_\lambda$ and does not depend on the asymptotic
constant in (SR.6).

\begin{proof}
\proofstep{1. Exact cancellations and the nonzero limit.}
Evaluation of (SR.1) at zero gives
$\int_0^\infty\phi_{j,\lambda}(u)du=\mu_jd_j/2$.
Consequently $\int_0^\infty\psi_\lambda=0$ exactly.
The uniform Hermite comparison alone gives
$\|\phi_{j,\lambda}-h_j\|_2=O(\lambda^{-3/2})$
after choosing the signs, including the omitted Gaussian
tail and the harmless $L^2$ renormalization.
Since $\mathcal Fh_j=h_j$ for $j=0,4$, it follows that
$1-\nu_j=\|(1-P_\lambda)\mathcal F\phi_{j,\lambda}\|_2^2
=O(\lambda^{-3})$ and $\mu_j\to1$.
Thus the uniform ratio comparison at zero implies
\[
 \sup_{|u|\leq\lambda}|\psi_\lambda(u)-h(u)|
                         \leq C\lambda^{-2},       \tag{SR.7}
\]
where $h$ is AC-S04.  To check the scalar, the Hermite
ratios at zero are
\[
 \frac{h_0(u)}{h_0(0)}=e^{-\pi u^2},\qquad
 \frac{h_4(u)}{h_4(0)}=
     \frac{16\pi^2u^4-24\pi u^2+3}{3}e^{-\pi u^2};
\]
their difference multiplied by $-3/4$ is $h(u)$.
The same comparisons give uniform upper bounds for
$\|\psi_\lambda\|_2$, $|d_0|^{-1}$ and $|d_4|^{-1}$.

For $\lambda^{-1}\leq u\leq\lambda$, summing (SR.7)
over $n\leq\lambda/u$ bounds the theta error by
$C\lambda^{-1}u^{-1/2}$.  Its squared $L^2(du/u)$
norm is at most
$C\lambda^{-2}\int_{\lambda^{-1}}^\lambda u^{-2}du
=O(\lambda^{-1})$.
The omitted Gaussian terms $nu>\lambda$ are negligible:
extracting $e^{-\pi\lambda^2/2}$ from
$|h(nu)|\leq C(1+(nu)^4)e^{-\pi(nu)^2}$ and bounding
the remaining Gaussian sum by $C(1+u^{-1})$ gives
an exponentially small bound after integration on this
window.  AC-S04 gives an exponentially small $L^2$
tail for $K$ outside the window.  Therefore
\[
              \|k_\lambda-K\|_2=O(\lambda^{-1/2}),
                                                        \tag{SR.8}
\]
and in particular $\|k_\lambda\|_2$ is bounded below.

\proofstep{2. A sharp Fourier tail in logarithmic Sobolev norms.}
Put $\phi=\mathcal F\psi_\lambda$ and
$\eta=(1-P_\lambda)\phi$.
Orthogonality of the prolate modes, both before and
after Fourier transformation, gives
\[
                    \|\eta\|_2^2\leq C\delta_\lambda.
                                                        \tag{SR.9}
\]
Indeed the exterior Fourier inner products are diagonal,
with diagonal entries $1-\nu_j$, and
$1-\nu_0\leq1-\nu_4$.

We need more than this $L^2$ estimate, because a theta
multiplier and a discrete sampling operation follow it.
The requisite coarse regularity has polynomial cost.
Differentiating (SR.1) inside the time interval and
using Plancherel shows
\[
 \|\phi_{j,\lambda}'\|_{L^2(-\lambda,\lambda)}
        \leq2\pi|\mu_j|^{-1}\|u\phi_{j,\lambda}\|_2
        \leq C\lambda.
\]
The same bound holds for $\psi_\lambda'$.
The elementary one-dimensional trace inequality then
bounds the endpoint values by $C\lambda^{1/2}$.
Consequently the total variations of the zero extensions
of $\psi_\lambda$ and $u\psi_\lambda$ are bounded by
$C\lambda^{3/2}$ and $C\lambda^{5/2}$ respectively.
Integration by parts in their Fourier transforms yields
\[
       |\phi(v)|\leq C\lambda^{3/2}|v|^{-1},
 \qquad |\phi'(v)|\leq C\lambda^{5/2}|v|^{-1}.
                                                        \tag{SR.10}
\]

Dilate the positive tail to endpoint one and then take
logarithms:
\[
       w(t)=1_{(0,\infty)}(t)
                   \sqrt\lambda e^{t/2}\phi(\lambda e^t).
                                                        \tag{SR.11}
\]
It satisfies $\|w\|_2\leq C\sqrt{\delta_\lambda}$ and,
for $t>0$,
\[
 |w(t)|\leq C\lambda^3e^{-t/2},\qquad
 |w'(t)|\leq C\lambda^3e^{t/2}.                    \tag{SR.12}
\]
For every nonnegative integer $m$ and $0<u<1/2$,
these estimates imply
\[
                  \|t^mw\|_{H^u(\mathbb R)}
                                  \leq C_{m,u}\lambda^3.
                                                        \tag{SR.13}
\]
Here the actual jump at $t=0$ is included.
For completeness, for $0<r<1$ split the translation
difference at $R=|\log r|$.
On $[0,R]$ the squared difference is bounded by $r^2$
times the integral of the squared derivative; on
$[R,\infty)$ use the squared amplitude of each translate.
The interval crossing zero has length $r$ and the same
bound.  Equation (SR.12), with the derivatives of $t^m$,
therefore gives
\[
 \|(t^mw)(\cdot+r)-t^mw\|_2^2
        \leq C_m\lambda^6 r(1+|\log r|)^{2m}.
\]
Its integral against $r^{-1-2u}dr$ is finite precisely
in the stated range; translations of size at least one
are bounded by the $L^2$ norm.  This proves (SR.13).

The small norm in (SR.9) can be retained with an
arbitrarily small loss.  For an integer $m>1$, H\"older
gives
\[
 \|tw\|_2\leq
       \|w\|_2^{1-1/m}\|t^mw\|_2^{1/m}.
\]
Interpolate this bound and the bound for $w$ with
(SR.13), using a fixed $u<1/2$.
Given $\varepsilon>0$, first take $m$ large and then
$\sigma>0$ sufficiently small.  Fourier H\"older
interpolation gives
\[
 \|w\|_{H^\sigma}+\|tw\|_{H^\sigma}
       \leq C_\varepsilon\lambda^{A_\varepsilon}
                    \delta_\lambda^{(1-\varepsilon)/2}.
                                                        \tag{SR.14}
\]
We will use this statement with a smaller loss than
the $\varepsilon$ finally prescribed in (SR.3).

\proofstep{3. The principal theta multiplier and its tails.}
In raw logarithmic coordinates $\mathcal Uf(t)=e^{t/2}f(e^t)$,
write $\mathcal E$ for the closed multiplier
$\zeta(1/2-i\tau)$.
For compactly supported functions of bounded variation,
it agrees with the corrected theta expression
\[
 E f(x)=\sqrt x\sum_{n\geq1}f(nx)
                      -x^{-1/2}\int_0^\infty f(v)dv.
                                                        \tag{SR.15}
\]
One direct verification uses
\[
 \sum_{n\leq Y}n^{-s}
       =\frac{Y^{1-s}}{1-s}+\zeta(s)+O_s(Y^{-\operatorname{Re}s})
                  \quad(0<\operatorname{Re}s<1).
\]
Integrate (SR.15) from $x=\epsilon$ to infinity
against $x^{s-1/2}dx/x$, interchange the finite sums,
and apply this identity with $Y=v/\epsilon$.
The displayed main term cancels the subtracted integral.
For bounded compact input the remaining error is
$O_s(\epsilon^{\operatorname{Re}s}(1+|\log\epsilon|))$.
Letting $\epsilon\downarrow0$ gives
$\zeta(s)\int_0^\infty f(v)v^{s-1}dv$.
This proves the multiplier identity on the critical line
whenever the resulting norm is finite, and then by
closure on its stated domain.

We record why $A$ makes its Sobolev order arbitrarily
small, including the derivative needed below:
\[
 |\zeta(1/2+i\tau)|+|\zeta'(1/2+i\tau)|
                 \leq C_b(1+|\tau|)^b
                         \quad\text{for every }b>0. \tag{SR.16}
\]
Apply the analytic logarithm argument of AC-T28 to
$(s-1)\zeta(s)/(s+1)$ in $\operatorname{Re}s>1/2$;
it is holomorphic and zero-free there under $A$.
The bounded Euler logarithm on disks centered at $2+i\tau$,
Borel--Caratheodory on a larger zero-free disk, and
three-circles give $|\log\zeta(1/2+d+i\tau)|
\leq C_d(\log(2+|\tau|))^{\beta_d}$, $\beta_d<1$,
up to a bounded rational-factor logarithm.
Thus the bound on that line is smaller than every
fixed positive power of $|\tau|$.
The functional equation gives on $1/2-d+i\tau$ the
same bound times $O((1+|\tau|)^d)$.
The polynomial-growth three-lines theorem, applied
after cancelling the pole, bounds the intervening
strip by an arbitrarily small power by choosing
$d$ small.  Cauchy's formula on a fixed smaller disk
in this strip proves the same assertion for $\zeta'$.
The remaining bounded range of $\tau$ is compact and
contains no pole.  This proves (SR.16).

Consequently, for arbitrarily small $s',b>0$,
\[
 \|\mathcal Ew\|_{H^{s'}}+
                 \|t\mathcal Ew\|_{H^{s'}}
       \leq C_{s',b}
          \bigl(\|w\|_{H^{s'+b}}+
                         \|tw\|_{H^{s'+b}}\bigr). \tag{SR.17}
\]
The second estimate follows by differentiating the
Mellin multiplier; its two terms use $\zeta$ and
$\zeta'$.  It also proves the asserted closed-domain
membership by Fourier multiplication.

There is no need to interpret a conditionally convergent
theta sum for the Fourier tail.  The cosine Mellin
identity for an even compact BV function, obtained by
Abel regularization of the cosine integral, gives
\[
 \int_0^\infty\phi(v)v^{s-1}dv
   =2(2\pi)^{-s}\Gamma(s)\cos(\pi s/2)
                     \int_0^\infty\psi_\lambda(u)u^{-s}du
                           \quad(0<\operatorname{Re}s<1).
\]
The regularization passes to the limit by the BV
$O(v^{-1})$ Fourier bound and the integrability of
$|\psi_\lambda(u)|u^{-\operatorname{Re}s}$ at zero.
Together with the zeta functional equation this proves
$E\psi_\lambda(x)=\mathcal E\phi(1/x)$ as an $L^2$
identity.  The compact part $P_\lambda\phi$ and the
tail $\eta$ belong to the multiplier domain, by
(SR.13), (SR.16), and the sharp-endpoint Sobolev
estimate for the compact prolate combination.
Linearity and (SR.15) therefore give on $x>\lambda$
\[
 \mathcal E\phi(x)=\mathcal E\eta(x)-c_Qx^{-1/2},
 \qquad
 c_Q=\int_0^\infty(P_\lambda\phi)(v)dv
       =-\frac38\mu_0\mu_4(\mu_0-\mu_4).
                                                        \tag{SR.18}
\]
Since $1-\nu_0\leq\delta_\lambda$ and
$1-\nu_4=\delta_\lambda$, $|c_Q|\leq C\delta_\lambda$.
This explicit scalar term is why no exact value-zero
constraint at the raw origin has to be added to the
source's two-mode trial.

After reflection and translation of logarithms, the
omitted lower theta tail is therefore represented by
\[
 R(t)=1_{(0,\infty)}(t)
       \left(\mathcal Ew(t)-c_Q\lambda^{-1/2}e^{-t/2}\right).
                                                        \tag{SR.19}
\]
Here $w$ in (SR.11) is exactly the logarithm of the
unitarily dilated raw tail $\sqrt\lambda\eta(\lambda u)$;
theta dilation covariance accounts for the factor
$\lambda^{-1/2}$ in the scalar term.

We also give the required sharp target restriction
estimate.  If $0<s<s'<1/2$, then
\[
           \|1_{(0,\infty)}v\|_{H^s}
                                     \leq C_{s,s'}\|v\|_{H^{s'}}.
                                                        \tag{SR.20}
\]
In the Gagliardo norm the extra cross-boundary term is
a constant times $\int_0^\infty|v(t)|^2t^{-2s}dt$.
For $t\geq1$ use $L^2$; on $(0,1)$ use the
embedding $H^{s'}(\mathbb R)\subset
L^{2/(1-2s')}(\mathbb R)$ and H\"older.
The remaining power is integrable because $s/s'<1$.
This proves (SR.20) including possible jumps.
It applies also to $tv$, since multiplication by $t$
commutes with the restriction.
Equations (SR.14), (SR.17)--(SR.20), with sufficiently
small $s',b$ and a sufficiently large moment index,
now give
\[
 \|R\|_{H^s}+\|tR\|_{H^s}
      \leq C_\varepsilon\lambda^{A_\varepsilon}
                     \delta_\lambda^{(1-\varepsilon)/2}.
                                                        \tag{SR.21}
\]
The elementary function $1_{t>0}e^{-t/2}$ and its
first moment belong to every $H^s$, $s<1/2$, so
its scalar term satisfies the same estimate.

\proofstep{4. Sampling at all zeros, including multiplicities.}
The zero ordinates obey
$\sum_{\gamma\in[j,j+1]}m_\gamma
\leq C\log(2+|j|)$.
This is the usual local consequence of the
Riemann--von Mangoldt estimate; it also follows from
the Jensen disk bound used in AC-T40, after cancelling
the zeta pole.  If $f,tf\in H^s$, $s>0$, then
\[
 \sum_\gamma m_\gamma|\widehat f(\gamma)|^2
    \leq C_s\bigl(\|f\|_{H^s}^2+\|tf\|_{H^s}^2\bigr).
                                                        \tag{SR.22}
\]
Indeed $f,tf\in L^2$ imply $f\in L^1$, so its
Fourier transform has its continuous representative.
On each unit interval the elementary $H^1$ pointwise
bound controls its supremum by the integral of
$|\widehat f|^2+|\widehat f'|^2$ over a slightly
larger interval.  Multiply by the local zero count,
sum intervals, and use
$\log(2+|\tau|)\leq C_s(1+\tau^2)^s$.
Plancherel and $\widehat f'=\widehat{-itf}$ prove
(SR.22).  No zero separation or simplicity is used.

The full theta function $F(x)=E\psi_\lambda(e^x)$
has zero Fourier transform at every critical zeta zero.
For a direct justification, its mean is zero and
the BV Riemann-sum estimate gives $F(x)=O(e^{x/2})$
at $-\infty$, while $F(x)=0$ for $x>L$.
Its Mellin identity, initially obtained by absolute
summation for $\operatorname{Re}s>1$, continues to
$\operatorname{Re}s>0$; the pole at $s=1$ is cancelled
by the exact raw mean zero.  Hence on the critical
line its transform is $\zeta(s)\int\psi_\lambda(u)u^{s-1}du$.

The compact function $k_\lambda$ is piecewise smooth
with finitely many jumps, so it belongs to $H^s$
for every $s<1/2$ and to the full closed Weil form
domain.  This last assertion also follows directly
from AC-T06: its logarithmic Fourier-symbol form is
bounded by every positive Sobolev norm, and its
remaining compressed prime and pole terms are bounded.
Smooth approximation of these finitely many jumps
therefore converges in form norm.
By AC-S03 and the vanishing of the transform of $F$,
\[
 q_L(k_\lambda)
       =\sum_\gamma m_\gamma|\widehat R(-\gamma)|^2.
                                                        \tag{SR.23}
\]
Reflection and translation introduce only the phase
$e^{i\gamma L}$, whose modulus is one.
Under $A$ all the ordinates are real; their multiset
is invariant under sign.  Apply (SR.22) and (SR.21),
choosing the small losses to total the prescribed
$\varepsilon$, to obtain (SR.3).

Finally (SR.8), the minimum principle and (SR.6) give
(SR.4), absorbing each fixed polynomial in $\lambda$
into an arbitrarily small loss in the Gaussian exponent.
The form is invariant under logarithmic reflection;
under $A$ its even and odd parts are nonnegative and
orthogonal for both form and Hilbert norm.
Thus the even part of $k_\lambda$ has no larger form
value, and it still converges to the nonzero even $K$.
This proves the final assertion.
\end{proof}

\begin{remark}[The remaining lowest-line comparison]
The upper exponential rate in (SR.4) strengthens AC-S05
and supplies a rigorous one-sided counterpart of the
comparison plotted in S25, Figure \texttt{fpro1}.
It does not give the plotted asymptotic equality.
For a normalized trial $v_\lambda$ and a simple bottom
eigenvector $u_\lambda$, AC-S07 still requires control of
\[
 \frac{q_L(v_\lambda)-a_1(L)}{a_2(L)-a_1(L)}.
\]
Equation (SR.3) controls the numerator from above;
neither a lower bound for this gap nor a coercive
comparison on the orthogonal complement has been
proved here.  Thus no convergence of the actual
lowest eigenline, and no resulting large-window
regularized determinant convergence, is asserted.
The construction is an analytic estimate in the
original Hilbert space, with no new motivic or
algebraic-geometric identification attached to it.
\end{remark}
